\chapter{Theoretical background} \label{chap:theoreticalbackground} 

\gls{ns} is a software development approach that prioritizes achieving software stability
through the use of standardized, modular components and interfaces. This theory is
informed by several scientific disciplines, including systems theory, mathematics, and
computer science, as well as some other software development approaches such as agile
development and domain-driven design.

The goal of this thesis is to investigate whether the philosophy of \gls{ca}
aligns with the goals of \gls{ns}. To do so, it is essential to have a
comprehensive understanding of software stability and the key concepts, principles, and
architectures that impact software stability.

This chapter begins by examining the concepts of software stability, evolvability, and
modularity, highlighting their significance in achieving software stability in \gls{ns}.
This is followed by a brief overview of the design theorems and proposed architecture of
\gls{ns}.

The subsequent sections of the thesis explore the fundamental principles that underlie
\gls{ca}, as well as its proposed architecture designs. Finally, the thesis
concludes by discussing which aspects of \gls{ca} align with the principles of
\gls{ns} and contribute to achieving software stability in this approach.

\section{Towards stable software architectures} \label{sec:on_stability}


\gls{ns} originated in the field of software engineering, aiming to achieve modular and
stable software artifacts. However, the underlying theory of \gls{ns} can be applied to
various other domains, such as Enterprise Engineering, Business Process Modeling, and
document management. This research acknowledges the software engineering background of
gls{ns}. It consistently refers to software and Information Systems when referring to
\enquote*{artifacts}. However, the reader should realize that the concepts and artifacts
are not restricted to software artifacts alone.

In several disciplines stability has been defined as \emph{Bounded Input Bounded Output}
(BIBO). It is the fundamental property of a system when subjected to bounded input
disturbances. BIBO stability ensures that the output of a system will also be bounded,
preventing uncontrolled or unexpected behavior \parencite[270]{mannaert_normalized_2016}. 

A real-world example of the importance of stability is the Tacoma Narrows Bridge in
Washington State, USA. The bridge, depicted in figure \ref*{fig:bridge}, collapsed on
November the 7th, 1940. This was caused due to wind-induced oscillations called
aeroelastic flutter. The wind (Input) induced oscillations in the bridge, causing it to
start swaying back and forth (Output). These oscillations were initially small, but as
they continued, they began to increase in amplitude or magnitude, causing the bridge to
collapse.

\begin{figure}[H]
    \centering
    \includegraphics[width=0.6\textwidth]{Figures/bridge.pdf}
    \caption[TNB]{Tacoma Narrows Bridge (Galloping Gertie)}
    \label{fig:bridge}
\end{figure}

Stability can also be used in the context of software engineering. In the context of
\gls{ns}, it is considered a critical property that ensures that the software is not
excessively sensitive to small changes \parencite[270]{mannaert_normalized_2016}. New
functional requirements should only lead to fixed, and an expected amount of changes in
the source code. Conversely, instabilities occur when the total number of modifications
relies on the size of the software artifact. The number of changes will grow over time in
parallel with the growth of the software artifact. These instabilities are referred to as
combinatorial effects \parencite[270]{mannaert_normalized_2016}. when combinatorial effects
are absent, the software artifact can be considered evolvable.
\section{Towards evolvable software architectures} \label{sec:on_modules}

There are a couple of aspects concerning evolvable software architectures. One of which is
stability, as described in chapter \ref{sec:on_stability} \nameref{sec:on_stability}.
Another factor that impacts evolvability is the modularity of the architecture. There is a
wide consensus about two fundamental rules when thinking of-, and designing modularity:
\emph{high cohesion} and \emph{low coupling} \autocite[22]{mannaert_normalized_2016}.
These will

Software modules can be defined as self-contained units of code that perform specific
tasks or sets of tasks within a larger system. A software module is designed to operate
independently of other modules, with well-defined interfaces that allow it to communicate
and exchange data with other modules if necessary \autocite[22]{mannaert_normalized_2016}.

A module can be considered a hierarchical and recursive concept. They are independent of
their size (lines of code) or computational magnitude. They can be as small as a function
as part of a class. The class itself can also be considered a module. A group of classes
contained in a Dynamic Link Library (DLL) or Application Programming Interface (API) can
also be considered a module of an even bigger system. 

An important part of the design of a software system is to identify the possible different
modules and their interaction interfaces.

\subsection{Cohesion: The beauty of Software Desing} \label{subsec:on_modularity}

The term cohesion denotes the extent to which the various structural components of a
software system operate cohesively towards a singular and well-defined objective or goal.
Empirical studies in software engineering have extensively demonstrated the significance
of cohesion, linking higher levels of cohesion with reduced defects, enhanced
maintainability, and greater openness to change. Consequently, achieving high cohesion has
been associated with an overall improvement in software quality attributes such as
reliability, maintainability, reusability, and evolvability.

Cohesion facilitates the reduction of complexity and interdependence among the components
of a system, thereby contributing to a more efficient, maintainable, and reliable system.
By organizing components around a shared purpose or function, or by standardizing their
interfaces, data structures, and protocols, cohesion can offer the following benefits:

\begin{itemize}
    \item \textbf{Reduce redundancy and duplication of effort}: \\
    Cohesion ensures that components are arranged around a common purpose or function,
    reducing duplicates or redundant code. This simplifies system comprehension,
    maintenance, and modification.
    \item \textbf{Promoting code reuse:}\\
    Cohesion facilitates code reuse by making it easier to extract and reuse components
    designed for specific functions. This saves time and effort during development and
    enhances overall system quality.
    \item \textbf{Enhance maintainability:}\\
    Cohesion decreases the complexity and interdependence of system components, making it
    easier to identify and rectify bugs or errors in the code. This improves system
    maintainability and reduces the risk of introducing new errors during maintenance.
    \item \textbf{Increase scalability:}\\
    Cohesion improves a system's scalability by enabling it to be extended or modified
    effortlessly to accommodate changing requirements or conditions. By designing
    well-organized and well-defined components, developers can easily add or modify
    functionality as needed without disrupting the rest of the system.  
\end{itemize}

\subsection{Coupling: The beast of Software Desing} \label{subsec:on_coupling}

Coupling is a central concept in software engineering that pertains to the degree of
interdependence among software modules or components. The level of coupling between
modules denotes the strength of their relationship, whereby a high level of coupling
implies a significant degree of interdependence. Conversely, low coupling signifies a
weaker relationship between modules, where modifications in one module are less likely to
impact others.

The negative impacts of excessive coupling on software systems are considerable. High
coupling can render software systems difficult to maintain, modify, or evolve. It can
impede the identification and resolution of errors within a system, leading to prolonged
debugging periods. Additionally, it can cause fragility in the system, where slight
modifications in one module can trigger cascading failures throughout the entire system.
Therefore, it is crucial for software engineers to minimize coupling between modules while
maintaining a cohesive design. By developing systems with low coupling, software engineers
can construct more maintainable, scalable, and adaptable systems that are easier to evolve
over time.

Coupling, in software engineering, can take several forms, including content, common,
control, stamp, and data coupling. Content coupling occurs when one module accesses or
modifies the internal data or logic of another module, leading to high interdependence and
difficulty in isolating errors. Common coupling occurs when several modules access and use
the same global data, increasing their interdependence and reducing modularity. Control
coupling occurs when one module controls the execution flow of another module, making it
difficult to modify or reuse the controlled module. Stamp coupling arises when two modules
share a common data structure, leading to tight coupling and high interdependence.
Finally, data coupling exists when two modules share data, which can lead to coupling
between them.

To avoid the negative impacts of coupling on software systems, software engineers should
aim to minimize the degree of coupling between modules. This can be achieved by designing
cohesive, loosely coupled modules with well-defined interfaces. Loose coupling enables
each module to operate independently, reducing the impact of modifications made to other
modules. By implementing modular design principles, such as high cohesion and low
coupling, software engineers can develop systems that are easier to maintain, test, and
evolve.

In conclusion, coupling is a critical concept in software engineering that can have a
considerable impact on the maintainability, flexibility, and scalability of software
systems. By minimizing coupling between modules, software engineers can develop more
robust, adaptable systems that are easier to modify and evolve. Adopting modular design
principles can also facilitate the development of cohesive, loosely coupled modules that
enable the independent operation and reduce the impact of modifications made to other
modules.
\input{Chapters/TheoreticalBackgroud/Sections/Combinatorics.tex}
\section{Cohesion: The Beauty of Software Design}
<<toelichten waarop Cohesion invloed heeft op de evolueerbaarheid van de software>>

Cohesion in software engineering refers to the degree to which the different structural
parts of a software system work together to achieve a single and well-defined purpose, a
common goal. 

There is a considerable body of scientific evidence supporting the importance of software
cohesion. Several studies correlate high levels of cohesion with fewer defects and are
likely to be more maintainable. It has shown that software engineering artifacts are more
open to change. Having a high cohesion is often referred to have a positive effect on
software quality attributes like reliability, maintainability, reusability and thus the
evolvability of the software artifact. 

\subsection{Coupling}
<<toelichten waarop Coupling invloed heeft op de evolueerbaarheid van de software>>
\subsection{The Design Theorems} \label{subsec:ns_desing_theorems}

\gls{ns} Theorems is a theoretical framework of principles that aims to enhance the
stability of a software artifact. \gls{ns} provides a rigorous mathematical foundation
that offers guidelines for designing and developing software systems. The principles of NS
have gained significant attention in both academic and industrial circles due to their
potential to increase the evolvability and stability of software artifacts while reducing
maintenance costs. The theorems have been scientifically established and proven. In the
following sections, we will focus on the theorems of \gls{ns}. When applicable, we
emphasize some manifestations in one of the artifacts to demonstrate the level of
convergence of \gls{ca}.
